\Cor
	Пусть $f \in \Q[x]$ "--- неприводимый многочлен степени $5$, имеющий над $\R$ ровно $3$ корня. Тогда уравнение $f(x) = 0$ не разрешимо в радикалах.

\Proof
	Можно считать, что каждый автоморфизм из $\Gal{f}$ осуществляет перестановку корней многочлена $f$ над $\overline{\Q}$, поэтому $\Gal{f} \subset S_5$. Пусть корни $\alpha_1, \alpha_2 \in \overline{\Q}$ многочлена $f$ "--- комплексные, а остальные "--- вещественные. Тогда комплексное сопряжение, то есть $(12) \in S_5$, лежит в $\Gal{f}$. В силу утверждения о продолжении гомоморфизма, для любых $i, j \in \{1, \dotsc, 5\}$ существует автоморфизм из $\Gal{f}$ такой, что $\alpha_i \mapsto \alpha_j$. Докажем теперь, что из этого следует равенство $\Gal{f} = S_5$.
	
	Рассмотрим граф транспозиций в $\Gal{f}$, то есть граф с множеством вершин $\{1, \dotsc, 5\}$ и множеством ребер $\{\{i, j\} : (ij) \in \Gal{f}\}$. Поскольку $\Gal{f}$ "--- группа, то компоненты связности в этом графе являются кликами. Заметим теперь, что при сопряжении перестановкой $\sigma \in \Gal{f}$ такой, что $\sigma(i) = j$, компонента, содержащая $i$, переходит в компоненту, содержащую $j$. Значит, все компоненты имеют одинаковый размер. Но число $5$ "--- простое, поэтому в графе есть лишь одна компонента, и, следовательно, $\Gal{f} = S_5$. Эта группа неразрешима, поэтому и уравнение $f(x) = 0$ не разрешимо в радикалах.
\EndProof

\Example
	Многочлен $f := x^5 - 10x + 5 \in \Q[x]$ не разрешим в радикалах. Действительно, исследуя $f$ на экстремумы, легко проверить, что он имеет ровно три корня над $\R$, а его неприводимость следует из признака Эйзенштейна по модулю $5$. Значит, применимо утверждение выше.